% !TeX root = ../despliegue.tex
% ===========================================================================
%  Requisito previo y arranque.
% ===========================================================================

\section{Requisito único}

\textbf{Docker 24 o superior, con Compose v2.} Nada más.

No hace falta instalar Java, Node, Maven ni PostgreSQL: todo se compila y se
ejecuta dentro de contenedores. Es una decisión deliberada —el Principio VII de
la constitución del proyecto exige que el sistema levante con un solo comando— y
también lo que hace que la evaluación no dependa de qué tenga instalado quien
evalúa.

Para comprobar la versión instalada:

\begin{comando}
docker --version        # Docker version 24.x o superior
docker compose version  # Docker Compose version v2.x
\end{comando}

\section{Levantar el sistema}

\begin{comando}
git clone <url-del-repositorio> && cd reserva-canchas
docker compose up
\end{comando}

Y abrir \url{http://localhost}.

El primer arranque compila las imágenes y tarda varios minutos; los siguientes
son cuestión de segundos, porque Docker reutiliza las capas ya construidas.

Para dejarlo corriendo en segundo plano y que el comando no devuelva el control
hasta que todo esté sano:

\begin{comando}
docker compose up -d --wait
\end{comando}

Cuando termina, \cmd{docker compose ps} debe mostrar los once servicios en
estado \cmd{healthy}.

\subsection{Si algún puerto está ocupado}

El gateway se publica en el puerto 8080 del anfitrión. Si esa máquina ya lo
usa, el arranque falla con \cmd{port is already allocated}. Se cambia sin tocar
el archivo de Compose, definiendo la variable al vuelo:

\begin{comando}
GATEWAY_PORT=18080 docker compose up -d --wait
\end{comando}

Todos los puertos publicados admiten esa sustitución:

\begin{comando}
SHELL_PORT              MS_USUARIOS_PORT        GATEWAY_PORT
MF_RESERVAS_PORT        MS_CANCHAS_PORT         POSTGRES_PORT
MF_ADMINISTRACION_PORT  MS_RESERVAS_PORT        EDGE_PORT
MF_REPORTES_PORT        MS_REPORTES_PORT
\end{comando}

\nota{Cambiar esos puertos no afecta a la aplicación. El navegador solo usa
\url{http://localhost}, que es el puerto 80 del \emph{edge}; los demás puertos
existen para inspeccionar cada servicio por separado.}
